% GENERATED FILE. Edit canonical lessons, then run npm run generate:books.
% canonical-source-hash: fnv1a64:a95ff10bbb4f30ad
% canonical-lessons: TA-C286-tonrukiratu, TA-C286-enave, TA-R286-opinion-and-result

\chapter{What I Think and So}
\label{ch:ta-what-i-think-and-so}
\hlchaptermodality{286}

\begin{chapteropening}
\textbf{By the end of this chapter:} \emph{I can give an opinion with … eṉṟu niṉaikkiṟēṉ, … eṉṟu tōṉṟukiṟatu and eṉ karuttil, and give a result with ataṉāl, ataṉāltāṉ and eṉavē.}

Complete the last lesson of chapter 286: give an opinion and say what followed, with tōṉṟukiṟatu, eṉ karuttil, ataṉāl and eṉavē.
\end{chapteropening}

\section[\texorpdfstring{… \ta{என்று} \ta{தோன்றுகிறது} (… eṉṟu tōṉṟukiṟatu)}{… என்று தோன்றுகிறது (… eṉṟu tōṉṟukiṟatu)}]{… \ta{என்று} \ta{தோன்றுகிறது} (eṉṟu tōṉṟukiṟatu) — it seems to me, in my opinion}

\label{lesson:TA-C286-tonrukiratu}



\begin{quote}
Say \textquotedblleft{}I think this is a good house.\textquotedblright{} (\emph{itu nalla vīṭu eṉṟu niṉaikkiṟēṉ}.) What closes the opinion before the verb? (\textbf{\ta{என்று}}.) Then say \textquotedblleft{}dampness\textquotedblright{}. (\emph{īram}.)
\end{quote}

\subsection*{You'll want to know: … \ta{என்று} \ta{தோன்றுகிறது}}
\begin{quote}
\textbf{\ta{எனக்கு} \ta{நாளை} \ta{மழை} \ta{பெய்யும்} \ta{என்று} \ta{தோன்றுகிறது}.}
\end{quote}

\begin{quote}
\emph{eṉakku nāḷai maḻai peyyum eṉṟu tōṉṟukiṟatu.} — \textquotedblleft{}It seems to me it will rain tomorrow.\textquotedblright{}
\end{quote}

\textbf{\ta{தோன்றுகிறது}} (\emph{tōṉṟukiṟatu}, \textquotedblleft{}it appears, it occurs\textquotedblright{}). The person it occurs to takes \textbf{-\ta{க்கு}}, as in \textbf{\ta{எனக்குத்} \ta{தமிழ்} \ta{தெரியும்}}: \textquotedblleft{}to me, '(it) will rain tomorrow', so it appears\textquotedblright{}. \textbf{\ta{பெய்யும்}} (\emph{peyyum}, \textquotedblleft{}it will fall\textquotedblright{}) is the \textquotedblleft{}it\textquotedblright{} form of the future.

\textbf{\ta{நினைக்கிறேன்}} says you have weighed it; \textbf{\ta{தோன்றுகிறது}} says it is an impression, and sounds softer. Both take the whole opinion first, closed by \textbf{\ta{என்று}}.

Two more shapes of the same \textbf{\ta{நினை}}: \textbf{… \ta{என்று} \ta{நினைத்தேன்}} (\emph{… eṉṟu niṉaittēṉ}, \textquotedblleft{}I thought that …\textquotedblright{}), and the negative, \textbf{\ta{அவன்} \ta{வருவான்} \ta{என்று} \ta{நான்} \ta{நினைக்கவில்லை}} (\emph{avaṉ varuvāṉ eṉṟu nāṉ niṉaikkavillai}, \textquotedblleft{}I don't think he will come\textquotedblright{}). \textbf{\ta{வருவான்}} (\emph{varuvāṉ}, \textquotedblleft{}he will come\textquotedblright{}).

\subsection*{You'll want to know: \ta{என்} \ta{கருத்தில்}, 'in my opinion'}
\begin{quote}
\textbf{\ta{என்} \ta{கருத்தில்}, \ta{ரயில்} \ta{பயணம்} \ta{நல்லது}.}
\end{quote}

\begin{quote}
\emph{eṉ karuttil, rayil payaṇam nallatu.} — \textquotedblleft{}In my opinion, travelling by train is good.\textquotedblright{}
\end{quote}

\textbf{\ta{கருத்து}} is \textquotedblleft{}an opinion\textquotedblright{}; \textbf{\ta{கருத்தில்}} is \textquotedblleft{}in (the) opinion\textquotedblright{}. \textbf{\ta{நல்லது}} (\emph{nallatu}, \textquotedblleft{}a good thing; good\textquotedblright{}) is \textbf{\ta{நல்ல}} made into a thing you can point at. \textbf{\ta{என்} \ta{கருத்துப்படி}} (\emph{eṉ karuttuppaṭi}, \textquotedblleft{}according to my opinion\textquotedblright{}) says the same with \textbf{-\ta{படி}}, \textquotedblleft{}according to\textquotedblright{}.

This one sits at the front, apart, so the listener knows that what follows is yours.

\subsection*{Your turn}
Say these aloud:

\begin{itemize}
\raggedright
  \item \textquotedblleft{}It seems to me it will rain tomorrow\textquotedblright{} — \emph{eṉakku nāḷai maḻai peyyum eṉṟu tōṉṟukiṟatu}
  \item \textquotedblleft{}I don't think he will come\textquotedblright{} — \emph{avaṉ varuvāṉ eṉṟu nāṉ niṉaikkavillai}
  \item \textquotedblleft{}In my opinion, travelling by train is good\textquotedblright{} — \emph{eṉ karuttil, rayil payaṇam nallatu}
  \item one opinion of your own about the weather, with \ta{தோன்றுகிறது}
\end{itemize}

\begin{tcolorbox}[breakable,colback=teal!4,colframe=teal!35!black,title={Before you move on}]
\textquotedblleft{}It seems to me that …\textquotedblright{}? (\textbf{\ta{எனக்கு} … \ta{என்று} \ta{தோன்றுகிறது}}.) Who takes \textbf{-\ta{க்கு}}? (\textbf{The person it occurs to.}) \textquotedblleft{}In my opinion\textquotedblright{}? (\textbf{\ta{என்} \ta{கருத்தில்}}.) \textquotedblleft{}I thought that …\textquotedblright{}? (\textbf{… \ta{என்று} \ta{நினைத்தேன்}}.)
\end{tcolorbox}
\section[\texorpdfstring{\ta{எனவே} (eṉavē)}{எனவே (eṉavē)}]{\ta{எனவே} (eṉavē) — so, that is why, therefore}

\label{lesson:TA-C286-enave}



\begin{quote}
Say \textquotedblleft{}so, because of that\textquotedblright{}. (\emph{ataṉāl}.) Then \textquotedblleft{}It seems to me it will rain tomorrow.\textquotedblright{} (\emph{eṉakku nāḷai maḻai peyyum eṉṟu tōṉṟukiṟatu}.)
\end{quote}

\subsection*{You'll want to know: \ta{அதனால்} and \ta{அதனால்தான்}}
\begin{quote}
\textbf{\ta{நேற்று} \ta{மழை} \ta{பெய்தது}. \ta{அதனால்} \ta{நான்} \ta{வீட்டில்} \ta{இருந்தேன்}.}
\end{quote}

\begin{quote}
\emph{nēṟṟu maḻai peytatu. ataṉāl nāṉ vīṭṭil iruntēṉ.} — \textquotedblleft{}It rained yesterday. So I stayed at home.\textquotedblright{}
\end{quote}

\textbf{\ta{அதனால்}} is \textbf{\ta{அதன்}}, the shape of \textbf{\ta{அது}} that carries endings, with \textbf{-\ta{ஆல்}}, \textquotedblleft{}by, because of\textquotedblright{}: \textquotedblleft{}because of that\textquotedblright{}. The cause comes first, in its own sentence; \textbf{\ta{அதனால்}} opens the result.

Add \textbf{-\ta{தான்}}, \textquotedblleft{}exactly, that very\textquotedblright{}, and it points back hard: \textquotedblleft{}that is why\textquotedblright{}.

\begin{quote}
\textbf{\ta{அவனுக்குக்} \ta{காய்ச்சல்}. \ta{அதனால்தான்} \ta{அவன்} \ta{வரவில்லை}.}
\end{quote}

\begin{quote}
\emph{avaṉukkuk kāyccal. ataṉāltāṉ avaṉ varavillai.} — \textquotedblleft{}He has a fever. That is why he did not come.\textquotedblright{}
\end{quote}

\subsection*{You'll want to know: \ta{எனவே} and \ta{ஆகையால்}, the written 'therefore'}
\begin{quote}
\textbf{\ta{பேருந்து} \ta{வரவில்லை}. \ta{எனவே} \ta{நாங்கள்} \ta{நடந்து} \ta{போனோம்}.}
\end{quote}

\begin{quote}
\emph{pēruntu varavillai. eṉavē nāṅkaḷ naṭantu pōṉōm.} — \textquotedblleft{}The bus did not come. Therefore we walked.\textquotedblright{}
\end{quote}

\textbf{\ta{எனவே}} and \textbf{\ta{ஆகையால்}} (\emph{ākaiyāl}, \textquotedblleft{}therefore\textquotedblright{}) do the job of \textbf{\ta{அதனால்}} in written Tamil: a notice, a letter, the news. They are the formal side of the two Tamils. In talk, \textbf{\ta{அதனால்}} carries the result.

\subsection*{Your turn}
Say these aloud:

\begin{itemize}
\raggedright
  \item \textquotedblleft{}It rained yesterday. So I stayed at home.\textquotedblright{} — \emph{nēṟṟu maḻai peytatu. ataṉāl nāṉ vīṭṭil iruntēṉ.}
  \item \textquotedblleft{}He has a fever. That is why he did not come.\textquotedblright{} — \emph{avaṉukkuk kāyccal. ataṉāltāṉ avaṉ varavillai.}
  \item \textquotedblleft{}The bus did not come. Therefore we walked.\textquotedblright{} — \emph{pēruntu varavillai. eṉavē nāṅkaḷ naṭantu pōṉōm.}
\end{itemize}

\begin{tcolorbox}[breakable,colback=teal!4,colframe=teal!35!black,title={Before you move on}]
\textquotedblleft{}So\textquotedblright{}? (\textbf{\ta{அதனால்}}.) \textquotedblleft{}That is why\textquotedblright{}? (\textbf{\ta{அதனால்தான்}}.) \textquotedblleft{}Therefore\textquotedblright{}, in writing? (\textbf{\ta{எனவே}}, or \textbf{\ta{ஆகையால்}}.) Which comes first, the cause or the result? (\textbf{The cause.})
\end{tcolorbox}
\section[(review)]{Review — opinion and result}

\label{lesson:TA-R286-opinion-and-result}



\begin{quote}
Say \textquotedblleft{}It seems to me it will rain tomorrow.\textquotedblright{} Then read \textbf{\ta{ஈரம்}} aloud: which letter opens it? (\textbf{\ta{ஈ}}, \emph{ī}.)
\end{quote}

\subsection*{Your turn}
Say these aloud:

\begin{itemize}
\raggedright
  \item In my opinion, this is a good shop.
  \item It seems to me he will come tomorrow.
  \item It rained. So I stayed at home.
  \item The bus did not come. Therefore we walked to school.
  \item one thing that happened to you, and what you did because of it, with \ta{அதனால்}
\end{itemize}

\begin{tcolorbox}[breakable,colback=teal!4,colframe=teal!35!black,title={Before you move on}]
\textquotedblleft{}It seems to me\textquotedblright{}? (\textbf{\ta{எனக்கு} … \ta{என்று} \ta{தோன்றுகிறது}}.) \textquotedblleft{}In my opinion\textquotedblright{}? (\textbf{\ta{என்} \ta{கருத்தில்}}.) \textquotedblleft{}That is why\textquotedblright{}? (\textbf{\ta{அதனால்தான்}}.) \textquotedblleft{}Therefore\textquotedblright{}, in writing? (\textbf{\ta{எனவே}}.)
\end{tcolorbox}
